\begin{bankssolutions}{Solutions to Implicit Exercises for Chapter 9}

\BanksImplicitSolution{I-CH09-001}
Raise the index in the consistency relation with the inverse metric.  The beta
function is
\[
  \beta^i=-G^{ij}\partial_jC,
\]
so it is the negative gradient of $C$ in the metric $G$.  Along an RG
trajectory,
\begin{align*}
  \frac{\dd C}{\dd t}
  &=\partial_iC\,\dot g^i
   =-G_{ij}\beta^i\beta^j\\
  &=-G^{ij}\partial_iC\,\partial_jC\leq0.
\end{align*}
Positive definiteness makes the inequality strict whenever $\beta\ne0$.
Thus $C$ is a Lyapunov function rather than a periodic coordinate on coupling
space.  In particular, a nonconstant periodic orbit would make $C$ decrease
and return to its original value, which is impossible.

Compactness of the forward orbit gives a lower bound for $C$, so
$C(g(t))\to C_\infty$ and
\[
  \int_0^\infty\dd t\,
  G_{ij}\beta^i\beta^j=C(g(0))-C_\infty<\infty.
\]
On the compact orbit the smooth function
$G_{ij}\beta^i\beta^j$ is uniformly continuous along the flow.  If it stayed
above a fixed positive number along an infinite sequence of time intervals,
the integral would diverge.  Hence $\beta(g(t))\to0$.  Equivalently, every
point in the omega-limit set has $\dot C=0$, and positive definiteness then
sets $\beta=0$ there.  An omega-limit set of a continuous trajectory is
connected.  If its fixed points are isolated, it contains one point $g_*$,
and therefore
\[
  \lim_{t\to\infty}g(t)=g_*.
\]

Write $g=g_*+u$ and let $H_{ij}=\partial_i\partial_jC|_*$.  To first order,
\[
  \dot u^i=M^i{}_ju^j,
  \qquad M^i{}_j=-G_*^{ik}H_{kj}.
\]
This matrix is self-adjoint in the $G_*$ inner product, so its eigenvalues are
real.  We orient $t$ toward the infrared, as in a step that lowers the cutoff.
An eigenmode $u_a\sim\ee^{y_at}$ grows for $y_a>0$ and is relevant.  A mode
with $y_a<0$ decays and is irrelevant; $y_a=0$ is marginal and requires
higher-order analysis.  For one coupling with $G=1$, $\dot g=-C'(g)$ and
$\dot C=-(C')^2$.  A local minimum of $C$ has
$\dot u=-C''(g_*)u$ and is infrared-attractive, which checks every sign in the
general formula.

\BanksImplicitSolution{I-CH09-002}
After integrating out the shell, the cutoff is $\ee^{-t}\Lambda$.  The change
of variables $x=\ee^t x'$ returns it to $\Lambda$.  A scaling field is chosen
so that
\[
  \Phi'(x')=\ee^{t\Delta_\Phi}\Phi(\ee^t x').
\]
At a fixed point the resulting dimensionless action equals the original one:
\[
  S_*[\Phi']=S_*[\Phi].
\]
This is precisely invariance under a dilatation, including the field
rescaling.  The anomalous part of $\Delta_\Phi$ records the rescaling needed
beyond its engineering dimension.

A scaling operator obeys
\[
  O_I(x)\longmapsto \ee^{-t\Delta_I}O_I(\ee^{-t}x).
\]
Since $\dd^D x\mapsto\ee^{Dt}\dd^D x$, its coupling changes according to
\[
  g^I(t)=\ee^{(D-\Delta_I)t}g^I(0)+O(g^2),
  \qquad \dot g^I=y_Ig^I+O(g^2),
  \qquad y_I=D-\Delta_I.
\]
With infrared-directed $t$, $y_I>0$, $y_I=0$, and $y_I<0$ give relevant,
marginal, and irrelevant perturbations, respectively.

Fixed-point invariance also fixes separated-point correlators.  For identical
scalar fields,
\[
  \left\langle\Phi(\lambda x_1)\cdots\Phi(\lambda x_n)\right\rangle_*
  =\lambda^{-n\Delta_\Phi}
  \left\langle\Phi(x_1)\cdots\Phi(x_n)\right\rangle_*.
\]
Spin indices acquire the corresponding finite-dimensional Lorentz matrices.

For a free massless scalar,
$S_*=(1/2)\int\dd^D x\,(\partial\phi)^2$ is invariant when
\[
  \Delta_\phi=\frac{D-2}{2}.
\]
Normal-ordered powers have
$\Delta_{\phi^n}=n(D-2)/2$, so their couplings have
$y_n=D-n(D-2)/2$.  In $D=4$, $\phi^2$ is relevant, $\phi^4$ is marginal,
and $\phi^n$ with $n>4$ is irrelevant.  This is the scaling classification
used throughout Banks's chapter.

\BanksImplicitSolution{I-CH09-003}
A standard gauge-fixing fermion has the schematic form
\[
  \Psi_\kappa=\int\dd^D x\,\bar c^a
  \left(\partial^\mu A^a_\mu-\kappa m_a\chi^a
        +\frac{\kappa}{2}B^a\right),
\]
where $\chi^a$ is a would-be Goldstone field and $B^a$ is auxiliary.  Its
precise normalization is immaterial for the argument.  Differentiation gives
\[
  \partial_\kappa S_\kappa=s(\partial_\kappa\Psi_\kappa).
\]
Let $\mathcal A_\kappa$ be a time-ordered LSZ amplitude between physical
states.  BRST invariance of the measure gives
\[
  \partial_\kappa\mathcal A_\kappa
  =\ii\left\langle\mathrm{out}\left|
  T\,s(\partial_\kappa\Psi_\kappa)\,
  \exp(\ii S_\kappa)\right|\mathrm{in}\right\rangle.
\]
Writing $sX=\ii[Q_{\mathrm B},X\}$ moves the BRST charge through the
time-ordered product.  It annihilates the BRST-invariant vacuum and both
cohomology classes,
\[
  Q_{\mathrm B}|\mathrm{in}\rangle=Q_{\mathrm B}|\mathrm{out}\rangle=0.
\]
The matrix element of the BRST-exact insertion therefore vanishes.  This step
uses an anomaly-free measure and the usual absence of boundary contributions
in field space.  It proves
\[
  \partial_\kappa\mathcal A_\kappa=0.
\]

In an $R_\kappa$ gauge the Goldstone and ghost masses are proportional to
$\sqrt{\kappa}\,m_a$.  Taking $\kappa\to\infty$ after renormalization, at fixed
physical masses and external data, removes them from finite-energy external
states and enforces $\chi^a=0$.  This is the unitary-gauge limit.  The preceding
identity makes its physical amplitude equal to the value at every finite
$\kappa$, provided this renormalized limit exists.  This order of limits is
important because the off-shell unitary-gauge power counting is singular.

The distinction between off-shell Green functions follows directly from the
propagators.  For one massive vector,
\[
  D^{(\kappa)}_{\mu\nu}(k)=
  \frac{-\ii}{k^2+m^2-\ii0}
  \left[\eta_{\mu\nu}-(1-\kappa)
  \frac{k_\mu k_\nu}{k^2+\kappa m^2-\ii0}\right]
\]
falls as $k^{-2}$ at fixed finite $\kappa$.  Its unitary-gauge limit is
\[
  D^U_{\mu\nu}(k)=
  \frac{-\ii}{k^2+m^2-\ii0}
  \left(\eta_{\mu\nu}+\frac{k_\mu k_\nu}{m^2}\right).
\]
The longitudinal numerator leaves an $O(m^{-2})$ term at large momentum.
Power counting of off-shell diagrams then permits counterterms of arbitrarily
high dimension.  These Green functions involve gauge-invariant, nonpolynomial
composites when rewritten in $R_\kappa$ variables, in accord with Banks's
explanation.

For a tree exchange between conserved currents, $k_\mu J^\mu=k_\nu J'^\nu=0$.
Every gauge-dependent longitudinal term drops out:
\[
  J^\mu D^{(\kappa)}_{\mu\nu}J'^\nu
  =J^\mu D^U_{\mu\nu}J'^\nu
  =\frac{-\ii J\mathbin{\cdot}J'}{k^2+m^2-\ii0}.
\]
This is the simplest explicit check of the cohomological proof.

\BanksImplicitSolution{I-CH09-004}
For a finite matrix $H_N(g)$, the exponential is defined by a power series
that converges uniformly on compact subsets of the complexified parameter
domain.  Each matrix entry of $\exp[-\beta H_N(g)]$ is therefore analytic, as
is its finite trace $Z_N$.  For real $\beta>0$ and Hermitian $H_N$,
\[
  Z_N=\sum_{r=1}^{\dim\mathcal H_N}\ee^{-\beta E_r}>0.
\]
At each real point, continuity gives a complex neighborhood free of zeros.
An analytic branch of $\log Z_N$ exists there, which proves real analyticity
of $F_N$ in $\beta$ and all real couplings.

For a finite-dimensional classical system,
\[
  Z_N(\beta,g)=\int_X\dd\mu(x)\,\ee^{-\beta H_N(x;g)}.
\]
Suppose the integrand extends holomorphically to a complex neighborhood of
each physical parameter point and has an integrable majorant there, uniform on
compact subsets.  Morera's theorem, with dominated convergence interchanging
the contour and $x$ integrals, makes $Z_N$ holomorphic.  Equivalently, Cauchy
estimates justify its locally convergent Taylor series term by term.  The same
positivity and logarithm argument then applies on the physical domain.

The thermodynamic free energy is a scaled limit
\[
  F(\beta,g)=\lim_{N\to\infty}
  -\frac{1}{\beta V_N}\log Z_N(\beta,g).
\]
Analyticity passes to a limit under locally uniform convergence of the
functions and the needed derivatives.  At a phase transition this uniformity
fails.  A sequence of smooth curves may then approach a function with a cusp,
or its derivatives may develop a jump or divergence.  In the equivalent
Lee--Yang description, every finite $Z_N$ has isolated zeros away from the
physical real domain.  As $V_N$ grows, their number grows and they can pinch a
real point.  No common zero-free neighborhood then survives, and
$V_N^{-1}\log Z_N$ can acquire a nonanalytic limit.

For one Ising spin $s=\pm1$ in a field $h$,
\[
  Z_1(\beta,h)=\sum_{s=\pm1}\ee^{\beta hs}=2\cosh(\beta h).
\]
It is positive for real $\beta h$ and its free energy is analytic there.  Its
zeros lie at $\beta h=\ii\pi(n+1/2)$, away from the real axis.  A phase
transition requires the infinite-volume accumulation mechanism described
above.

\BanksImplicitSolution{I-CH09-005}
The exact blocked action is defined by
\[
  \ee^{-H_1[\Phi]}=
  \sum_{\{s_x\}}\ee^{-H[s]}
  \prod_I\delta\!\left(
  |B_I|\Phi_I-\sum_{x\in B_I}s_x\right).
\]
Split the constrained microscopic variables into the retained block field and
fast variables $\chi$, and write their coupling as $V(\Phi,\chi)$.  Then
\[
  H_1[\Phi]=H_{\mathrm{slow}}[\Phi]
  -\log\left\langle\ee^{-V(\Phi,\chi)}\right\rangle_{\!\mathrm{fast}}
  =H_{\mathrm{slow}}[\Phi]
  +\sum_{r\geq1}\frac{(-1)^{r+1}}{r!}
  \left\langle V(\Phi,\chi)^r\right\rangle_{\!c,\mathrm{fast}}.
\]
Expanding each $V$ in $\Phi$ expresses every coefficient $K_n$ as a sum of
connected fast-mode cumulants.  This is the precise linked-cluster meaning of
the kernel expansion; conventional factorials depend on the definition of
$K_n$.

By assumption, a connected fast-mode cumulant is bounded by an exponential
of the length of a shortest tree joining its blocks:
\[
  |K_n(I_1,\ldots,I_n)|
  \leq C_n\exp[-c\,\ell_{\mathrm{tree}}(I_1,\ldots,I_n)].
\]
This is the linked-cluster statement: disconnected clusters disappear from
the logarithm, while a connected cluster reaching widely separated blocks
must contain a chain of short-range microscopic interactions across that
distance.  The induced interactions are therefore quasi-local on the blocked
lattice.  Each later step integrates another band of short-distance modes.
Convolution of exponentially decaying kernels remains exponentially decaying,
with a changed range measured in the new lattice units.  Iteration preserves
quasi-locality.

If a block after $k$ steps contains $M=B^{kd}$ Ising spins, then
\[
  \Phi=-1,-1+\frac2M,\ldots,1-\frac2M,1.
\]
Its mesh $2/M$ tends to zero.  After the field-strength rescaling appropriate
to the fixed point, sums over these values become Riemann sums and then a
functional integral over a continuous field.  The microscopic bound can be
encoded by a local potential that grows rapidly outside the image of
$[-1,1]$; long-distance correlators are insensitive to the detailed shape of
that wall.

For independent unbiased spins,
\[
  \Pr(\Phi=m)=2^{-M}{M\choose M(1+m)/2}.
\]
Stirling's formula gives, for $m=O(M^{-1/2})$,
\[
  \Pr(\Phi=m)\simeq
  \sqrt{\frac{2}{\pi M}}\,
  \exp\!\left(-\frac{Mm^2}{2}\right).
\]
The rescaled variable $\varphi=\sqrt M\,\Phi$ has spacing $2/\sqrt M$ and
tends to a continuous unit Gaussian.  This explicit limit realizes Banks's
replacement of a finely spaced block variable by a continuum field.

\BanksImplicitSolution{I-CH09-006}
For an internal line of momentum $q$, the assumptions imply a bound of the
form
\[
  \left|\frac1{K(q^2/\Lambda^2)}\right|
  \leq C(1+|q|)^a\exp[-c q^2/\Lambda^2]
\]
outside a fixed ball.  A vertex contributes at most another polynomial.  Route
the internal momenta of a graph as
\[
  q_e=A_{er}\ell_r+P_e,
\]
where $\ell_r$, $r=1,\ldots,L$, are independent loop momenta and $P_e$ is a
linear combination of external momenta.  The cycle-edge matrix has rank $L$.
Consequently, for constants depending on the graph,
\[
  \sum_e|q_e|^2\geq c_G\sum_{r=1}^L|\ell_r|^2-C_G(P).
\]
The absolute value of the graph integrand is therefore bounded by
\[
  C_P(1+|\ell|)^N
  \exp\!\left[-c'_G\sum_r|\ell_r|^2/\Lambda^2\right].
\]
This function is integrable on $\mathbb R^{DL}$.  Applying the same routing
argument with the complementary loop momenta fixed proves convergence of
every UV subintegration, so the estimate also excludes subdivergences.

At a fixed monomial in the couplings and a fixed number of external legs,
there are finitely many graphs.  Each is finite by the bound above.  This is
the coefficient-by-coefficient meaning of a finite perturbation series; it
does not assert convergence of the infinite sum over perturbative orders.

A massless propagator can still produce a singularity at small momentum.
Move a positive quadratic term from $g_2$ into the free kernel,
\[
  K(p)\longmapsto K(p)+m^2,
  \qquad S_I\longmapsto S_I-\frac12m^2\int\phi^2.
\]
The full action stays fixed.  The new propagator is bounded at $p=0$, while
its large-$p$ exponential decrease is unchanged.  The infrared rearrangement
therefore leaves the ultraviolet proof intact.

\BanksImplicitSolution{I-CH09-007}
For a Weyl variation, the defining metric formula gives
\[
  \delta_\sigma W[g]
  =\int\dd^D x\sqrt g\,\sigma(x)
  \langle T^\mu{}_{\mu}(x)\rangle.
\]
Insert any collection $X$ of renormalized local operators at points outside
the support of $\sigma$.  The same Ward identity reads
\[
  \int\dd^D x\sqrt g\,\sigma(x)
  \langle T^\mu{}_{\mu}(x)X\rangle=0.
\]
It holds for every compactly supported test function $\sigma$.  The
fundamental lemma for distributions therefore implies
\[
  \langle T^\mu{}_{\mu}(x)X\rangle=0.
\]
Since $X$ is arbitrary, $T^\mu{}_{\mu}=0$ as an operator-valued distribution,
up to the contact terms already covered by the renormalized Ward identity.
The assumptions about anomalies and the virial current are exactly what
allows this representative of the stress tensor.

A constant rescaling alone would give only
$\int\sqrt g\,T^\mu{}_{\mu}=0$ and would not support the distributional step.
The Ward identity for arbitrary local test functions, supplied by the stated
no-anomaly and no-virial assumptions, is essential here.

Let $g_{\mu\nu}(t)=\ee^{2t\sigma(x)}g_{\mu\nu}$.  Along this path,
\[
  \frac{\dd}{\dd t}W[g(t)]
  =\int\dd^D x\sqrt{g(t)}\,\sigma(x)
  \langle T^\mu{}_{\mu}\rangle_{g(t)}=0.
\]
Integration from $t=0$ to $1$ proves invariance under the corresponding finite
Weyl transformation.

In flat space a conformal Killing vector satisfies
\[
  \partial_\mu\xi_\nu+\partial_\nu\xi_\mu=2\omega(x)\eta_{\mu\nu}.
\]
The diffeomorphism generated by $\xi$ changes the metric by this expression.
A Weyl rescaling with parameter $-\omega$ restores $\eta_{\mu\nu}$.  Diffeomorphism
covariance together with Weyl invariance therefore gives the flat-space
conformal Ward identity.

For a free massless scalar in $D>2$, use the improved tensor
\[
  T_{\mu\nu}=\partial_\mu\phi\partial_\nu\phi
  -\frac12\eta_{\mu\nu}(\partial\phi)^2
  +\xi_c(\eta_{\mu\nu}\Box-\partial_\mu\partial_\nu)\phi^2,
  \qquad
  \xi_c=\frac{D-2}{4(D-1)}.
\]
Taking the trace and using $\Box\phi=0$ gives
$T^\mu{}_{\mu}=0$.  This checks both the improvement coefficient and the
local Weyl conclusion.

\BanksImplicitSolution{I-CH09-008}
Expanding $\exp(-\delta S)$ in the Euclidean functional integral gives a
two-leg insertion $-m^2$.  A line with $r$ insertions contributes
\[
  G_0(p)[-m^2G_0(p)]^r.
\]
Summing over $r$ gives
\begin{align*}
  G(p)&=\frac1{p^2}\sum_{r=0}^\infty
  \left(-\frac{m^2}{p^2}\right)^r\\
  &=\frac1{p^2}\frac1{1+m^2/p^2}
   =\frac1{p^2+m^2}.
\end{align*}
The ordinary geometric series converges for $|m^2|<|p^2|$.  Elsewhere the
last expression defines its analytic continuation, or equivalently the exact
Gaussian integral with the quadratic term included in the inverse
propagator.

With mostly-plus Minkowski action
$\mathcal L=-\frac12(\partial\phi)^2-\frac12m^2\phi^2$, the free massless line is
$-\ii/(p^2-\ii0)$ and the insertion is $-\ii m^2$.  Hence
\[
  \frac{-\ii}{p^2-\ii0}
  \sum_{r=0}^\infty
  \left[\frac{-\ii m^2(-\ii)}{p^2-\ii0}\right]^r
  =\frac{-\ii}{p^2+m^2-\ii0}.
\]
The Euclidean result tends to $1/p^2$ when $m^2\to0$.  It also obeys the
Dyson equation
\[
  G=G_0-G_0m^2G,
\]
whose solution is $(G_0^{-1}+m^2)^{-1}$.  This is why quadratic relevant
operators are resummed before treating the remaining CFT perturbations.

\BanksImplicitSolution{I-CH09-009}
Put $L=16\pi^2$ and $d_{\mathrm{reg}}=4-\epsilon$.  The one-loop scalar bubble in any channel is
\begin{align*}
  B(P^2)&=\mu^\epsilon\int\frac{\dd^{d_{\mathrm{reg}}}k}{(2\pi)^{d_{\mathrm{reg}}}}
  \frac{1}{(k^2+m^2)[(k+P)^2+m^2]}\\
  &=\frac{\mu^\epsilon}{(4\pi)^{d_{\mathrm{reg}}/2}}
  \Gamma\!\left(2-\frac {d_{\mathrm{reg}}}{2}\right)
  \int_0^1\dd x\,[m^2+x(1-x)P^2]^{d_{\mathrm{reg}}/2-2}\\
  &=\frac{2}{L\epsilon}+B_{\mathrm{fin}}(P^2,m^2)+O(\epsilon).
\end{align*}
Each of the $s,t,u$ graphs has symmetry factor $1/2$.  Their common pole is
therefore
\[
  3\left(\frac{\lambda^2}{2}\right)
  \frac{2}{L\epsilon}=\frac{3\lambda^2}{L\epsilon}.
\]
This gives the one-loop term in the bare coupling displayed in the exercise.

Ignoring self-energy insertions for the moment, there are two two-loop 1PI
skeletons.  With labelled external legs, the iterated-bubble family has three
members, one in each channel, and symmetry factor $1/4$.  The overlapping
family has six members: for each channel there is a wineglass and its inverse,
each with symmetry factor $1/2$.  A representative, with all external momenta
incoming and $P=p_1+p_2$, is
\[
  \begin{aligned}
    W(p_i)&=\mu^{2\epsilon}
    \int\frac{\dd^{d_{\mathrm{reg}}}k\,\dd^{d_{\mathrm{reg}}}\ell}
    {(2\pi)^{2d_{\mathrm{reg}}}}\\[-2pt]
    &\quad\times
    \frac{1}{D(k)D(P-k)D(\ell)D(k+p_3-\ell)},
    \qquad D(r)=r^2+m^2.
  \end{aligned}
\]
The iterated graph is the corresponding product of two $B$ integrals.
Permuting the external legs in $W$ produces all six overlapping graphs.
There is also a bubble with a one-loop tadpole inserted in an internal line.
Its tadpole subdivergence is canceled by the one-loop mass counterterm; in a
massless calculation both terms are scaleless.  It supplies no new four-point
counterterm.

Schwinger parameters put either skeleton in the form
\[
  \frac{\mu^{2\epsilon}}{(4\pi)^{d_{\mathrm{reg}}}}
  \int_{x_i\geq0}\![\dd x]\,
  \frac{\Gamma(\epsilon)\,
  \mathcal P(x)}{\mathcal U(x)^{d_{\mathrm{reg}}/2}}
  [\mathcal M^2(x,p)]^{-\epsilon},
\]
with the appropriate graph polynomials $\mathcal U$, $\mathcal P$, and
$\mathcal M^2$.  A boundary where a one-loop four-point subgraph shrinks gives
a pole multiplying the finite part of the reduced one-loop graph.  There are
six insertions obtained by replacing either vertex of each one-loop channel by
\[
  \delta\lambda^{(1)}=\frac{3\lambda^2}{L\epsilon}.
\]
For each forest, this insertion cancels the nonlocal
$B_{\mathrm{fin}}(P^2)/\epsilon$ term of its parent graph.  The tadpole forests
cancel against the mass insertion described above.  The remaining parameter
integrals are finite at nonexceptional momenta, and every pole left after the
forest subtractions is independent of the external momenta.

To display the normalization cleanly, let $Z_4$ multiply the renormalized
$\lambda\phi^4/4!$ vertex and let
$\phi_0=Z_\phi^{1/2}\phi$.  Summing the three iterated graphs, the six
wineglass graphs, and their one-loop counterterm insertions gives
\[
  Z_4=1+\frac{3\lambda}{L\epsilon}
  +\frac{\lambda^2}{L^2}
  \left(\frac9{\epsilon^2}-\frac3\epsilon\right)
  +O(\lambda^3).
\]
The two-loop self-energy gives
\[
  Z_\phi=1-\frac{\lambda^2}{12L^2\epsilon}+O(\lambda^3).
\]
Thus the external-field subtraction changes the simple two-loop coefficient
through
\[
  \lambda_0=\mu^\epsilon\lambda Z_4Z_\phi^{-2}.
\]
It converts $-3$ into $-3+1/6=-17/6$, while leaving the double pole unchanged.
Consequently
\[
  \lambda_0=\mu^\epsilon\left[
  \lambda+\frac{3\lambda^2}{L\epsilon}
  +\frac{\lambda^3}{L^2}
  \left(\frac9{\epsilon^2}-\frac{17}{6\epsilon}\right)
  \right]+O(\lambda^4).
\]
The momentum independence of its poles is the locality check: a single local
$\phi^4$ counterterm removes the full four-point divergence.

The double pole supplies a useful internal check.  Write
\[
  \lambda_0=\mu^\epsilon
  \left(\lambda+\frac{a_1(\lambda)}\epsilon
  +\frac{a_2(\lambda)}{\epsilon^2}+\cdots\right).
\]
Finiteness of $\mu\dd\lambda/\dd\mu$ requires at order $\lambda^3$
$a_2=9\lambda^3/L^2$, the square dictated by the coefficient
$3\lambda^2/L$ of the one-loop subdivergence.  It is therefore fixed before
the overall two-loop integral is evaluated.

Differentiating the bare relation at fixed $\lambda_0$ gives
\[
  \beta(\lambda)=-\epsilon\lambda+\lambda a_1'(\lambda)-a_1(\lambda).
\]
Here
\[
  a_1(\lambda)=\frac{3\lambda^2}{L}
  -\frac{17\lambda^3}{6L^2},
\]
and hence
\[
  \beta(\lambda)=-\epsilon\lambda+\frac{3\lambda^2}{L}
  -\frac{17\lambda^3}{3L^2}+O(\lambda^4).
\]
Together with the two-loop two-point subtractions for
$(\partial\phi)^2$ and $m^2\phi^2$ computed in Banks's discussion, this
four-point subtraction closes the program on the three operators
$(\partial\phi)^2$, $m^2\phi^2$, and $\phi^4$.  Setting $\lambda=0$ recovers
the free theory, while the momentum independence of every pole checks
locality.

\BanksImplicitSolution{I-CH09-010}
Set $d_{\mathrm{reg}}=4-\epsilon$ for the Euclidean loop integrals.
Let
\[
  b_f=\frac43\sum_fq_f^2,
  \qquad b_s=\frac13\sum_sq_s^2,
  \qquad b=b_f+b_s.
\]
Flavor traces simply add because an Abelian photon vertex preserves the
species circulating in a loop.  A Dirac loop gives
\[
  \Pi^{(f)}_{\mu\nu}\big|_{\mathrm{div}}
  =-\frac{e^2q_f^2}{6\pi^2\epsilon}
  (p^2\delta_{\mu\nu}-p_\mu p_\nu).
\]
For a complex scalar, the bubble and seagull combine as
\[
  \Pi^{(s)}_{\mu\nu}=e^2q_s^2\int\frac{\dd^{d_{\mathrm{reg}}}k}{(2\pi)^{d_{\mathrm{reg}}}}
  \left[
  \frac{(2k+p)_\mu(2k+p)_\nu}
  {(k^2+m_s^2)[(k+p)^2+m_s^2]}
  -\frac{2\delta_{\mu\nu}}{k^2+m_s^2}
  \right].
\]
Contracting with $p^\mu$ changes the first numerator into the difference of
the two scalar inverse propagators.  A shift of $k$ cancels that difference
against the seagull.  Thus the result is transverse, and its pole is
\[
  \Pi^{(s)}_{\mu\nu}\big|_{\mathrm{div}}
  =-\frac{e^2q_s^2}{24\pi^2\epsilon}
  (p^2\delta_{\mu\nu}-p_\mu p_\nu).
\]
Summing species gives
\[
  Z_3=1-\frac{e^2}{8\pi^2\epsilon}
  \left(\frac43\sum_fq_f^2+\frac13\sum_sq_s^2\right)+O(e^4).
\]

The Abelian Ward identity fixes
\[
  e_0=\mu^{\epsilon/2}eZ_3^{-1/2}.
\]
Differentiating at fixed $e_0$ and keeping the leading pole yields
\[
  \beta(e)=-\frac\epsilon2e+\frac{b e^3}{16\pi^2}+O(e^5).
\]
Its four-dimensional value is the formula in the exercise.  Every term in
$b$ is positive, which gives the infrared freedom emphasized by Banks.

The fermion calculation differs between species only through $q_f^2$ and the
mass.  Its divergent self-energy is
\[
  \Sigma_f^\infty(p)=-\frac{e^2q_f^2}{8\pi^2\epsilon}
  \left[\kappa(\slashed p-\ii m_f)+3\ii m_f\right].
\]
It is canceled by
\[
  Z_{2f}=1+\frac{\kappa e^2q_f^2}{8\pi^2\epsilon},
  \qquad
  \delta m_f=-\frac{3e^2q_f^2}{8\pi^2\epsilon}m_f.
\]
The Ward identity sets the pole part of the fermion-photon vertex equal to
that of $Z_{2f}$, so no independent charge counterterm appears.

For a scalar, photon exchange gives the pole
\[
  \Sigma_s^\infty(p)=-C_s
  \left[(3-\kappa)(p^2+m_s^2)-3m_s^2\right],
  \qquad C_s=\frac{e^2q_s^2}{8\pi^2\epsilon}.
\]
In the parametrization
$Z_s[p^2+m_s^2+\delta m_s^2]$, it is removed by
\[
  Z_s=1+C_s(3-\kappa),
  \qquad \delta m_s^2=-3C_s m_s^2.
\]
These are the gauge contributions.  The scalar potential also contributes to
the mass counterterm.  To keep every charge-allowed quartic without choosing a
special flavor basis, resolve the complex fields into real components and
write
\[
  V=\frac12(M^2)_{AB}\varphi_A\varphi_B
  +\frac1{4!}\lambda_{ABCD}\varphi_A\varphi_B\varphi_C\varphi_D,
\]
where gauge invariance restricts the symmetric tensor $\lambda_{ABCD}$.  Its
one-loop tadpole is
\[
  \begin{aligned}
    \Sigma^{(\lambda)}_{AB}
    &=\frac12\lambda_{ABCD}\mu^\epsilon
    \int\frac{\dd^{d_{\mathrm{reg}}}k}{(2\pi)^{d_{\mathrm{reg}}}}
    \left[(k^2\mathbf 1+M^2)^{-1}\right]_{CD},\\[-2pt]
    \Sigma^{(\lambda)}_{AB}\big|_{\mathrm{div}}
    &=-\frac{\lambda_{ABCD}(M^2)_{CD}}{16\pi^2\epsilon}.
  \end{aligned}
\]
It is momentum independent and is canceled by the opposite local mass-matrix
counterterm.  The massless photon seagull tadpole is scaleless.  Thus the full
scalar two-point subtraction is the sum of this potential term and the gauge
terms displayed above.  A basis that diagonalizes $M^2$ within each
equal-charge sector recovers the species-by-species formula.
The three-point scalar-photon Ward identity fixes its vertex pole from $Z_s$.
The four-point seagull identity then fixes the $A^2\phi_s^*\phi_s$ pole.  The
gauge dependence stays in field normalizations, while the displayed mass
shift is gauge-parameter independent in this parametrization.

Scalar and photon boxes also generate local four-scalar poles.  Their charge
flow permits exactly the monomials allowed by gauge invariance.  For generic
charges these include $|\phi_s|^2|\phi_t|^2$; special charge relations can
allow more quartics, scalar cubics, and Yukawa interactions.  They must be
included unless an explicit flavor or $\phi_s\mapsto-\phi_s$ symmetry forbids
them.

Here is a normalization check that separates the proper vertex from the field
counterterm.  For one unit-charge scalar with canonical kinetic term and
$V=\lambda|\phi|^4/4$, in Landau gauge the one-loop factors are
\[
  Z_{\lambda,\mathrm{vert}}\lambda
  =\lambda+\frac{5\lambda^2+24e^4}{16\pi^2\epsilon},
  \qquad
  Z_s=1+\frac{6e^2}{16\pi^2\epsilon}.
\]
Therefore
\[
  \lambda_0=\mu^\epsilon
  Z_s^{-2}Z_{\lambda,\mathrm{vert}}\lambda
  =\mu^\epsilon\left[
  \lambda+\frac{5\lambda^2-12e^2\lambda+24e^4}
  {16\pi^2\epsilon}\right]+O(\hbar^2).
\]
This also shows where the mixed field-normalization term enters.  With several
fields, the scalar, photon, and mixed diagrams contract the full quartic tensor
and the charges in every allowed channel.

Power counting restricts the divergent one-loop local terms to
$F^2$, the matter kinetic and mass operators, allowed Yukawa terms, and the
allowed scalar potential through fourth order.  BRST or Abelian Ward identities
tie their gauge vertices to the corresponding kinetic factors.  This proves
closure of the general renormalizable Lagrangian.  A single unit-charge Dirac
field gives $b=4/3$ and reproduces Banks's result.  One unit-charge complex
scalar gives $b=1/3$, exactly one quarter of the Dirac contribution.

\BanksImplicitSolution{I-CH09-011}
It is clearest to use the inverse-propagator form of the Ward identity.  With
the Euclidean factors absorbed consistently into $S^{-1}$,
\[
  p^\alpha\gamma_\alpha=S^{-1}(q+p)-S^{-1}(q).
\]
Hence
\begin{align*}
  p^\alpha\Pi_{\alpha\beta}(p)
  &=e_0^2\int\frac{\dd^{d_{\mathrm{reg}}}q}{(2\pi)^{d_{\mathrm{reg}}}}
  \operatorname{tr}\!\left[
  \{S^{-1}(q+p)-S^{-1}(q)\}
  S(q)\gamma_\beta S(q+p)\right]\\
  &=e_0^2\int\frac{\dd^{d_{\mathrm{reg}}}q}{(2\pi)^{d_{\mathrm{reg}}}}
  \operatorname{tr}\!\left[\gamma_\beta S(q)
  -\gamma_\beta S(q+p)\right].
\end{align*}
Cyclicity of the trace was used in the second line.  The factor
$\gamma_\beta$ is essential in both terms.  Dimensional regularization defines
the two analytically continued integrals with the same translation-invariant
measure.  Shifting $q\mapsto q-p$ in the second therefore gives
\[
  p^\alpha\Pi_{\alpha\beta}(p)=0.
\]

A sharp cutoff region $|q|<\Lambda$ moves under this shift and can leave a
surface term.  Gauge-invariant regularization, or counterterms fixed by the
Ward identity, removes that regulator artifact.  At $p=0$ the two terms agree
before integration.  Euclidean rotational invariance then permits
\[
  \Pi_{\alpha\beta}=A(p^2)\delta_{\alpha\beta}
  +B(p^2)p_\alpha p_\beta.
\]
Transversality sets $A=-p^2B$, so
\[
  \Pi_{\alpha\beta}(p)=
  (p^2\delta_{\alpha\beta}-p_\alpha p_\beta)\Pi(p^2).
\]
This excludes a local photon-mass counterterm.

\BanksImplicitSolution{I-CH09-012}
Integration by parts gives
\begin{align*}
  \Gamma(t+1)&=\int_0^\infty\dd u\,u^t\ee^{-u}\\
  &=\left[-u^t\ee^{-u}\right]_0^\infty
    +t\int_0^\infty\dd u\,u^{t-1}\ee^{-u}
   =t\Gamma(t).
\end{align*}
The boundary term vanishes for $\Re t>0$.  Since $\Gamma(1)=1$, iteration
gives $\Gamma(n)=(n-1)!$.

The beta integral follows by changing variables in a product of two gamma
integrals:
\[
  \Gamma(a)\Gamma(b)=\Gamma(a+b)
  \int_0^1\dd x\,x^{a-1}(1-x)^{b-1}.
\]
Taking $a=t$, $b=1-t$, and putting $x=y/(1+y)$ gives
\[
  \Gamma(t)\Gamma(1-t)
  =\int_0^\infty\dd y\,\frac{y^{t-1}}{1+y}.
\]
A keyhole contour for $z^{t-1}/(1+z)$ has one pole at $z=-1$ and yields
\[
  \int_0^\infty\dd y\,\frac{y^{t-1}}{1+y}
  =\frac{\pi}{\sin\pi t},
  \qquad 0<\Re t<1.
\]
This is Euler's reflection formula.  At $t=1/2$ it gives
$\Gamma(1/2)^2=\pi$.  The defining integral is positive, so the required root
is $\Gamma(1/2)=\sqrt\pi$.

For $M^2>0$, set $u=sM^2$:
\[
  \int_0^\infty\dd s\,s^{a-1}\ee^{-sM^2}
  =\Gamma(a)(M^2)^{-a}.
\]
At $a=1$ this is $1/M^2$, which is Schwinger's representation of an inverse
quadratic form.  The one-dimensional Gaussian is
$\int_{-\infty}^\infty\dd q\,\ee^{-sq^2}=\sqrt{\pi/s}$, so multiplication of
$d_{\mathrm{reg}}$ components gives
\[
  \int\dd^{d_{\mathrm{reg}}}q\,\ee^{-sq^2}
  =\left(\frac\pi s\right)^{d_{\mathrm{reg}}/2}.
\]
Dividing by $(2\pi)^{d_{\mathrm{reg}}}$ gives
$(2\sqrt\pi)^{-d_{\mathrm{reg}}}s^{-d_{\mathrm{reg}}/2}$, the factor in Banks's calculation.  These
formulae first hold in their convergence domains.  Their gamma-function
expressions provide the meromorphic continuation in $a$ and $d_{\mathrm{reg}}$ used by
dimensional regularization.

\BanksImplicitSolution{I-CH09-013}
The integrand is exactly a derivative:
\begin{align*}
  d_{\mathrm{reg}}\frac{\dd}{\dd s}
  \left(s^{1-d_{\mathrm{reg}}/2}\ee^{-sm_0^2}\right)
  =s^{1-d_{\mathrm{reg}}/2}\ee^{-sm_0^2}
  \left[-d_{\mathrm{reg}}m_0^2
  +\frac{d_{\mathrm{reg}}(2-d_{\mathrm{reg}})}{2s}\right].
\end{align*}
For $m_0^2>0$ and $\Re d_{\mathrm{reg}}<2$, the quantity in parentheses vanishes at both
$s=0$ and $s=\infty$.  The integral is therefore zero in this nonempty
convergence strip.  Both sides are meromorphic functions of $d_{\mathrm{reg}}$, so analytic
continuation gives the same identity near $d_{\mathrm{reg}}=4$.

This proves that the right-hand side of the equation for
$(d_{\mathrm{reg}}-1)p^2\Pi(p^2)$ vanishes at $p^2=0$.  It does not set $\Pi(0)$ to zero.
Indeed, $\Pi(0)$ contains the logarithmic pole that renormalizes $Z_3$.
Instead, the result says that the full tensor has no constant term:
\[
  \Pi_{\mu\nu}(p)=
  (p^2\delta_{\mu\nu}-p_\mu p_\nu)\Pi(p^2),
\]
with no $1/p^2$ singularity in $\Pi$ for positive fermion mass.  Thus the
photon pole remains at zero mass.

The $s^{1-d_{\mathrm{reg}}/2}$ factor is forced by the Jacobian
$\dd\alpha\,\dd\beta=s\,\dd s\,\dd x$ together with the Gaussian factor
$s^{-d_{\mathrm{reg}}/2}$.  Retaining it is what turns the zero-momentum expression into the
total derivative above.  For a massless fermion, keep $m_0>0$ as an infrared
regulator, establish the dimensionally continued identity, and then take
$m_0\to0$.  Setting $m_0=0$ before controlling the large-$s$ boundary can
mix this ultraviolet statement with the known infrared exceptional cases.

\BanksImplicitSolution{I-CH09-014}
Expand the momentum-dependent exponential:
\[
  \ee^{-sx(1-x)p^2}
  =\sum_{n=0}^\infty\frac{[-x(1-x)p^2]^n}{n!}s^n.
\]
The coefficient of $(p^2)^n$ in $\Pi$ is therefore proportional near the UV
endpoint to
\[
  \int_0\dd s\,s^{1-d_{\mathrm{reg}}/2+n}.
\]
At $d_{\mathrm{reg}}=4$ this is $\int_0\dd s\,s^{n-1}$.  The $n=0$ term is logarithmically
divergent.  Every $n\geq1$ term is integrable at $s=0$.  Factors of
$\ee^{-sm_0^2}$ affect the large-$s$ endpoint and leave this UV count
unchanged.

Multiplication by the transverse tensor supplies two powers of momentum:
\[
  \Pi_{\mu\nu}=(p^2\delta_{\mu\nu}-p_\mu p_\nu)
  [c_0+c_1p^2+c_2(p^2)^2+\cdots].
\]
Only the order-$p^2$ tensor term proportional to $c_0$ can diverge.  All terms
of order $p^4$ and higher are UV finite.  Direct differentiation confirms the
same point:
\[
  \left.\frac{\partial\Pi}{\partial p^2}\right|_{p^2=0}
  =-C_{d_{\mathrm{reg}}}\int_0^\infty\dd s\,s^{2-d_{\mathrm{reg}}/2}\int_0^1\dd x\,
  x(1-x)A_{d_{\mathrm{reg}}}(x)\ee^{-sm_0^2},
\]
whose small-$s$ power is $s^0$ at $d_{\mathrm{reg}}=4$.

\BanksImplicitSolution{I-CH09-015}
Through first order in $\alpha$, write the inverse propagator as
\[
  \Gamma_R^{(2)}(p)=
  (\slashed p-\ii m_0)
  +(Z_2-1)(\slashed p-\ii m_0)
  -\ii\delta m+\Sigma^\infty(p)+\Gamma_{\mathrm{fin}}^{(2)}(p).
\]
The field counterterm is
\[
  (Z_2-1)(\slashed p-\ii m_0)
  =\frac{\kappa\alpha}{2\pi\epsilon}
  (\slashed p-\ii m_0),
\]
which cancels the first term in $\Sigma^\infty$.  The mass counterterm gives
\[
  -\ii\delta m=+\frac{3\ii\alpha}{2\pi\epsilon}m_0,
\]
and cancels its remaining pole.  Replacing $m_0$ by $m_0+\delta m$ inside the
$Z_2-1$ term would produce $(Z_2-1)\delta m=O(\alpha^2)$, beyond the order of
this calculation.

The relation
\[
  m_0+\delta m=m_0
  \left(1-\frac{3\alpha}{2\pi\epsilon}\right)
\]
is multiplicative.  Its pole coefficient contains no $\kappa$.  At $m_0=0$
the mass counterterm vanishes, as required by chiral symmetry.  At
$\kappa=0$ the wave-function pole vanishes, while the same mass subtraction
remains.  These limits separately check both pieces of Banks's formula.

\BanksImplicitSolution{I-CH09-016}
Choose the infrared scale $Q=\sqrt{|\Delta|}$.  Along an RG characteristic the
two-point equation integrates to
\[
  F(\Delta/\mu^2,\alpha,\kappa)=
  \exp\!\left[-2\int_Q^\mu\frac{\dd k}{k}\,
  \gamma_F(\alpha(k),\kappa(k))\right]F(1,\alpha(Q),\kappa(Q)).
\]
The matching factor at $Q$ is regular and can be set to one at leading-log
accuracy.  Since $\beta=O(\alpha^2)$ while $\gamma_F=O(\alpha)$, the first
leading-log approximation holds $\alpha$ and $\kappa$ fixed.  It gives
\[
  F=\left(\frac{\Delta}{\mu^2}\right)^{\gamma_F^{(1)}},
  \qquad
  S_R(p)\simeq
  (\slashed p+\ii m)\mu^{-2\gamma_F^{(1)}}
  \Delta^{-1+\gamma_F^{(1)}}.
\]
Using Banks's one-loop value,
\[
  S_R(p)\ \propto\
  \frac{\slashed p+\ii m}
  {\Delta^{\,1+\kappa\alpha/(4\pi)}}.
\]
A noninteger exponent turns the isolated pole into a branch point and resums
the series of powers $[\alpha\ln(\Delta/\mu^2)]^n$.

Running of the infrared-free coupling gives
\[
  \alpha(k)=\frac{\alpha(\mu)}
  {1+\dfrac{2\alpha(\mu)}{3\pi}\ln(\mu/k)}
  +O(\alpha^2).
\]
If a scale-dependent gauge convention keeps $\kappa$ fixed while this charge
running is inserted, the characteristic factor becomes
\[
  F=\left[\frac{\alpha(\mu)}{\alpha(Q)}\right]^{3\kappa/4}.
\]
For the RG derivative at fixed bare gauge parameter, use instead
$\kappa_0=Z_3\kappa$.  Together with the Ward identity for the charge, this
gives in four dimensions
\[
  \frac{\dd\kappa}{\dd\ln k}
  =-\frac{2\alpha}{3\pi}\kappa+O(\alpha^2\kappa),
  \qquad
  \alpha(k)\kappa(k)=\alpha(\mu)\kappa(\mu).
\]
Thus $\gamma_F^{(1)}=-\kappa\alpha/(4\pi)$ is constant along the coupled
one-loop characteristic, and the running-value integral reproduces the pure
noninteger power above.  Higher-loop RG functions supply its logarithmic
corrections.

For the mostly-plus Minkowski continuation, set
$x=-(p^2+m^2)>0$ on the physical side of the cut.  Then
continuation gives
$\Delta\to-x\mp\ii0$.  Thus
\[
  (-x-\ii0)^{-1+\gamma_F}
  -(-x+\ii0)^{-1+\gamma_F}
  =2\ii\sin(\pi\gamma_F)x^{-1+\gamma_F},
\]
up to the chosen orientation of the discontinuity.  This is nonzero for the
perturbative noninteger exponent.  In Landau gauge, $\kappa=0$ and
$\gamma_F^{(1)}=0$, so the leading branch exponent returns to a simple pole.
Its gauge dependence shows that the elementary fermion field is not a
physical charged interpolating operator.  A gauge-invariant Coulomb dressing
and the coherent soft-photon states needed for inclusive scattering complete
the infrared construction.

\BanksImplicitSolution{I-CH09-017}
At $L$ loops, a primitive logarithmic divergence supplies one simple pole.
Nested or disjoint divergent subgraphs multiply their pole factors.  A graph
with $L$ independent loops contains at most $L$ such factors, so the forest
formula permits
\[
  \frac1\epsilon,\frac1{\epsilon^2},\ldots,\frac1{\epsilon^L}
\]
and no higher pole.  The highest pole is a product of lower-loop subtraction
factors selected by a maximal nested or disjoint forest.

For $r=1$, bare-coupling independence gives
\[
  0=\mu\frac{\dd g_0}{\dd\mu}
  =\mu^\epsilon\left[\epsilon F+\beta(g,\epsilon)\partial_gF\right].
\]
Write $\beta=-\epsilon g+\beta_4(g)$.  The finite term in the Laurent
expansion is
\[
  a_1-ga_1'+\beta_4=0,
\]
and hence
\[
  \boxed{\ \beta_4(g)=g a_1'(g)-a_1(g)\ }.
\]
The coefficient of $1/\epsilon^n$ gives, for $n\geq1$,
\[
  g a_{n+1}'-a_{n+1}=\beta_4 a_n'.
\]
These pole equations determine every higher-pole coefficient recursively from
lower data.  They are the algebraic form of the cancellation Banks describes.

For a field factor,
\begin{align*}
  \gamma_\Phi
  &=\frac12(-\epsilon g+\beta_4)
  \frac{\partial}{\partial g}
  \sum_{n\geq1}\frac{z_n(g)}{\epsilon^n}\\
  &=-\frac12g z_1'(g)
  +\frac12\sum_{n\geq1}\frac{
  -g z_{n+1}'(g)+\beta_4z_n'(g)}{\epsilon^n}.
\end{align*}
Finiteness imposes
\[
  g z_{n+1}'=\beta_4z_n',
\]
so
\[
  \boxed{\ \gamma_\Phi(g)=-\frac12g z_1'(g)\ }.
\]
Thus the residue of the simple pole supplies the anomalous dimension, while
the multiple poles enforce consistency.

For several couplings with
$g_0^a=\mu^{r_a\epsilon}[g^a+a_1^a/\epsilon+\cdots]$, the same calculation
gives
\[
  \beta_4^a=\sum_b r_b g^b\partial_ba_1^a-r_a a_1^a,
  \qquad
  \gamma_\Phi=-\frac12\sum_b r_bg^b\partial_bz_1.
\]
The higher-pole equations are, for $n\geq1$,
\[
  \left(\sum_b r_bg^b\partial_b-r_a\right)a_{n+1}^a
  =\sum_b\beta_4^b\partial_ba_n^a,
  \qquad
  \sum_b r_bg^b\partial_bz_{n+1}
  =\sum_b\beta_4^b\partial_bz_n.
\]
At one loop the forest has one element and only a simple pole occurs.  At two
loops one nested one-loop subgraph can multiply the overall pole, producing a
double pole; the recursion fixes it from the one-loop residue.  These checks
match the explicit $\phi^4$ result in Exercise~\ref{exercise:I-CH09-009}.

\BanksImplicitSolution{I-CH09-018}
At a fixed point with $\alpha_*\ne0$,
\[
  0=\beta(\alpha_*)=\alpha_*\gamma_A(\alpha_*),
\]
so $\gamma_A(\alpha_*)=0$.  Gauge invariance ties the field-strength
renormalization to that of the photon, and therefore
\[
  \Delta_F=2+\gamma_A=2.
\]
In spinor notation,
\[
  F_{\mu\nu}\longleftrightarrow
  (F_{\alpha\beta},\bar F_{\dot\alpha\dot\beta}),
\]
where the two symmetric spinors transform in $(1,0)$ and $(0,1)$.

Radial quantization turns a primary into a state.  For a primary in $(j,0)$
with $j>0$, positivity of the level-one descendant gives the conformal
unitarity bound
\[
  \Delta\geq j+1.
\]
For $F_{\alpha\beta}$, $j=1$, so the bound is saturated at $\Delta=2$.
The norm of the descendant
$P^{\alpha\dot\alpha}|F_{\alpha\beta}\rangle$ is proportional to
$\Delta-2$ and hence vanishes.  Hilbert-space positivity identifies a
zero-norm state with zero.  The operator-state correspondence then gives
\[
  \partial^{\alpha\dot\alpha}F_{\alpha\beta}=0.
\]
The conjugate representation similarly obeys
\[
  \partial^{\alpha\dot\alpha}\bar F_{\dot\alpha\dot\beta}=0.
\]
Together these are the source-free Maxwell equations and Bianchi identity for
the two helicities.

Conformal covariance fixes the two-point function of this primary to the free
Maxwell form, up to its positive normalization.  The null equation itself is
an operator statement: the norm of
$\partial^{\alpha\dot\alpha}F_{\alpha\beta}|0\rangle$ vanishes, positivity
removes the state, and the operator-state correspondence removes the local
operator.  Under the standard Wightman assumptions, the Jost--Schroer
free-field theorem then applies to the resulting local field with its free
two-point function.  Its commutator is the Maxwell c-number commutator and the
vacuum representation of the algebra generated by $F$ is the Maxwell Fock
representation.  The complementary local algebra supplies the multiplicity
space, so the $F$ algebra is a free tensor factor and Wick's theorem holds in
that factor.

Introduce a source $J^{\mu\nu}$ for $F_{\mu\nu}$.  The free-field theorem,
rather than the linear equation by itself, makes the connected generating
functional Gaussian:
\[
  W_F[J]=\frac12\int J\,\langle FF\rangle\,J.
\]
All derivatives of $W_F$ above second order vanish, so
\[
  \langle F_1\cdots F_n\rangle_{\!c}=0,
  \qquad n>2.
\]
The argument assumes unitarity, four-dimensional conformal invariance, a
gauge-invariant primary $F$ with nonzero two-point function, locality, the
spectrum condition, a unique cyclic vacuum, and a resolved primary in any
operator-mixing sector.  These are the hypotheses behind the free-sector
conclusion.

\BanksImplicitSolution{I-CH09-019}
Use
\[
  A_\mu=A^3_\mu,\qquad
  W^+_\mu=\frac{A^1_\mu-\ii A^2_\mu}{\sqrt2},\qquad
  W^-_\mu=(W^+_\mu)^*,
\]
and define
\[
  f_{\mu\nu}=\partial_\mu A_\nu-\partial_\nu A_\mu,\qquad
  Q_{\mu\nu}=W^+_\mu W^-_\nu-W^+_\nu W^-_\mu.
\]
With $D_\mu W^+_\nu=(\partial_\mu-\ii eA_\mu)W^+_\nu$, direct substitution
gives
\[
  F^3_{\mu\nu}=f_{\mu\nu}-\ii eQ_{\mu\nu},\qquad
  F^+_{\mu\nu}=D_\mu W^+_\nu-D_\nu W^+_\mu,\qquad
  F^-_{\mu\nu}=(F^+_{\mu\nu})^*.
\]
Since
$F^aF^a=F^3F^3+2F^+F^-$, the Yang--Mills Lagrangian becomes
\begin{align*}
  \mathcal L_{\mathrm{YM}}={}&-\frac14f_{\mu\nu}f^{\mu\nu}
  -\frac12
  (D_\mu W^+_\nu-D_\nu W^+_\mu)
  (D^\mu W^{-\nu}-D^\nu W^{-\mu})\\
  &+\ii e f_{\mu\nu}W^{+\mu}W^{-\nu}
  +\frac{e^2}{4}Q_{\mu\nu}Q^{\mu\nu}.
\end{align*}
The cross term follows from
$f^{\mu\nu}Q_{\mu\nu}=2f^{\mu\nu}W^+_\mu W^-_\nu$.
This also fixes the quartic coefficient, since both arise from the same
square $(F^3)^2$.

For comparison, a charged Proca field can have
\[
  \mathcal L_W=-\frac12|D_{[\mu}W_{\nu]}|^2-m^2W^+_\mu W^{-\mu}
  +\ii e\kappa_m f_{\mu\nu}W^{+\mu}W^{-\nu}.
\]
Minimal covariantization corresponds to $\kappa_m=0$.  Commuting two covariant
derivatives in its field equation supplies one unit of magnetic coupling; the
explicit Pauli term supplies $\kappa_m$ more.  For slowly varying fields, the
spatial equation gives
\[
  H_{\mathrm{mag}}=
  -\frac{e}{2m}(1+\kappa_m)\,
  \mathbf S\mathbin{\cdot}\mathbf B,
  \qquad (S_k)_{ij}=-\ii\epsilon_{kij}.
\]
Comparison with
$H_{\mathrm{mag}}=-(g_{\mathrm{mag}}e/2m)\mathbf S\cdot\mathbf B$
therefore gives
\[
  g_{\mathrm{mag}}=1+\kappa_m.
\]
The Yang--Mills expansion has $\kappa_m=1$, so
$g_{\mathrm{mag}}=2$.  Sending $e\to0$ removes the magnetic and quartic
terms and leaves the free vector kinetic term, which checks the decomposition.

\BanksImplicitSolution{I-CH09-020}
Set $d_{\mathrm{reg}}=4-\epsilon$ for the dimensionally regulated loop
integrals.
Write every pole in the form
\[
  \Pi_{\mu\nu}{}^{ab}\big|_{\mathrm{div}}
  =\frac{\ii g^2\delta^{ab}}{16\pi^2\epsilon}
  \left(Aq^2\eta_{\mu\nu}+Bq_\mu q_\nu\right).
\]
Terms proportional to a momentum-independent $\eta_{\mu\nu}$ are scaleless in
the massless theory.  Dimensional regularization sets them to zero.

For illustration, the ghost numerator is
$-(p+q)_\mu p_\nu$.  Combining denominators, shifting
$r=p+xq$, and retaining even terms gives
\[
  -(p+q)_\mu p_\nu
  \longmapsto-r_\mu r_\nu+x(1-x)q_\mu q_\nu.
\]
The pole integrals are
\[
  \int\frac{\dd^{d_{\mathrm{reg}}}r}{(2\pi)^{d_{\mathrm{reg}}}}
  \frac1{(r^2+\Delta)^2}
  =\frac{2}{16\pi^2\epsilon},\qquad
  \int\frac{\dd^{d_{\mathrm{reg}}}r}{(2\pi)^{d_{\mathrm{reg}}}}
  \frac{r_\mu r_\nu}{(r^2+\Delta)^2}
  =-\frac{\Delta\,\delta_{\mu\nu}}{16\pi^2\epsilon},
\]
with $\Delta=x(1-x)q_E^2$.  Since
$\int_0^1x(1-x)\dd x=1/6$, this gives the ghost row below.  Contracting the
two three-gluon vertices in Feynman gauge gives
$(A,B)=(19/6,-11/3)C_2(G)$.  The longitudinal part of either propagator is
  $-(1-\kappa)k_\rho k_\sigma/k^4$.  After symmetric integration its two possible
insertions add $(1-\kappa)(1,-1)C_2(G)$, while the term with two longitudinal
propagators cancels.  Thus the general-gauge three-gluon row is
$(25/6-\kappa,\kappa-14/3)C_2(G)$.  The four-gluon tadpole is scaleless.  The
components are therefore
\[
\begin{array}{c|cc}
  \text{graph}&A/C_2(G)&B/C_2(G)\\ \hline
  \text{two three-gluon vertices}&25/6-\kappa&\kappa-14/3\\
  \text{four-gluon tadpole}&0&0\\
  \text{ghost loop}&1/6&1/3\\ \hline
  \text{pure gauge plus ghost}&13/3-\kappa&-(13/3-\kappa)
\end{array}
\]
The first and third rows each have
$A+B\ne0$.  Their sum has $A+B=0$, which is the one-loop Slavnov--Taylor
identity.  Thus
\[
  \Pi_{\mu\nu}{}^{ab}\big|_{\mathrm{div,pure}}
  =\frac{\ii g^2C_2(G)\delta^{ab}}{16\pi^2\epsilon}
  \left(\frac{13}{3}-\kappa\right)
  (q^2\eta_{\mu\nu}-q_\mu q_\nu).
\]

Matter loops are separately transverse after the complex-scalar bubble is
combined with its seagull.  For collections of Weyl fermions and complex
scalars,
\[
\begin{array}{c|cc}
  \text{matter}&A&B\\ \hline
  \text{Weyl fermions}&-\frac43D(R_F)&+\frac43D(R_F)\\
  \text{complex scalars}&-\frac23D(R_S)&+\frac23D(R_S)
\end{array}
\]
in the same common normalization.  A Dirac fermion counts as two Weyl
fermions, and a real scalar in a real representation counts as one half of a
complex scalar.  The complete pole is therefore
\[
\begin{aligned}
  \Pi_{\mu\nu}{}^{ab}\big|_{\mathrm{div}}
  ={}&\frac{\ii g^2\delta^{ab}}{16\pi^2\epsilon}
  (q^2\eta_{\mu\nu}-q_\mu q_\nu)\\
  &\times\left[
  C_2(G)\left(\frac{13}{3}-\kappa\right)
  -\frac43D(R_F)-\frac23D(R_S)\right].
\end{aligned}
\]
This gives every tensor component because Lorentz covariance permits only the
two tensors tabulated above.

In Feynman gauge, $\kappa=1$, the pure coefficient is $10C_2(G)/3$.
Its trace is $(d_{\mathrm{reg}}-1)q^2$ times that coefficient and agrees with Banks's
component-contracted calculation.  This coefficient renormalizes the
covariant-gauge gluon field.  The coupling beta function follows from a
Slavnov--Taylor-related vertex or a gauge-invariant definition such as the
background-field effective action.  Its pure coefficient is $11C_2(G)/3$,
rather than the gauge-dependent $13C_2(G)/3-\kappa$.  This distinction is the
final consistency check.

\BanksImplicitSolution{I-CH09-021}
Use the Schwinger formula
\[
  \frac1{(p_E^2+U)^n}
  =\frac1{\Gamma(n)}\int_0^\infty\dd s\,
  s^{n-1}\ee^{-s(p_E^2+U)}.
\]
The Gaussian integrals are
\[
  \int\frac{\dd^{d_{\mathrm{reg}}}p_E}{(2\pi)^{d_{\mathrm{reg}}}}
  \ee^{-sp_E^2}=(4\pi s)^{-d_{\mathrm{reg}}/2},
\]
and
\[
  \int\frac{\dd^{d_{\mathrm{reg}}}p_E}{(2\pi)^{d_{\mathrm{reg}}}}
  p_{E\mu}p_{E\nu}\ee^{-sp_E^2}
  =\frac{\delta_{\mu\nu}}{2s}(4\pi s)^{-d_{\mathrm{reg}}/2}.
\]
The proper-time integral then gives
\[
  I_n^E(U)=\frac1{(4\pi)^{d_{\mathrm{reg}}/2}}
  \frac{\Gamma(n-d_{\mathrm{reg}}/2)}{\Gamma(n)}
  U^{d_{\mathrm{reg}}/2-n},
\]
and
\[
  I_{n,\mu\nu}^E(U)=
  \frac{\delta_{\mu\nu}}{2(4\pi)^{d_{\mathrm{reg}}/2}}
  \frac{\Gamma(n-d_{\mathrm{reg}}/2-1)}{\Gamma(n)}
  U^{d_{\mathrm{reg}}/2+1-n}.
\]
They first converge for $\Re n>\Re d_{\mathrm{reg}}/2$ and
$\Re n>\Re d_{\mathrm{reg}}/2+1$, respectively.  The right-hand sides define their
dimensional continuations.

Under the Feynman Wick rotation, $\dd^{d_{\mathrm{reg}}}p=\ii\dd^{d_{\mathrm{reg}}}p_E$ and
$p^2+U-\ii0=-(p_E^2+U)$.  Tensor components rotate into the Lorentz-covariant
factor $\eta^{\mu\nu}$.  Hence
\[
  \boxed{\displaystyle
  \int\frac{\dd^{d_{\mathrm{reg}}}p}{(2\pi)^{d_{\mathrm{reg}}}}
  \frac1{(p^2+U-\ii0)^n}
  =\frac{(-1)^n\ii}{2^{d_{\mathrm{reg}}}\pi^{d_{\mathrm{reg}}/2}}
  \frac{\Gamma(n-d_{\mathrm{reg}}/2)}{\Gamma(n)}
  \left(\frac1U\right)^{n-d_{\mathrm{reg}}/2}},
\]
and
\[
  \boxed{\displaystyle
  \int\frac{\dd^{d_{\mathrm{reg}}}p}{(2\pi)^{d_{\mathrm{reg}}}}
  \frac{p^\mu p^\nu}{(p^2+U-\ii0)^n}
  =\frac{(-1)^{n-1}\ii\eta^{\mu\nu}}
  {2^{d_{\mathrm{reg}}+1}\pi^{d_{\mathrm{reg}}/2}}
  \frac{\Gamma(n-d_{\mathrm{reg}}/2-1)}{\Gamma(n)}
  \left(\frac1U\right)^{n-d_{\mathrm{reg}}/2-1}}.
\]
The metric tensor is required by Lorentz covariance.

The scalar integral has mass dimension $d_{\mathrm{reg}}-2n$, so its power is
$U^{d_{\mathrm{reg}}/2-n}$.  In the tensor integral the numerator adds two
dimensions, giving $U^{d_{\mathrm{reg}}/2+1-n}$.  These exponents are the
negatives of the gamma-function arguments $n-d_{\mathrm{reg}}/2$ and
$n-d_{\mathrm{reg}}/2-1$.  This is the correspondence Banks asks the
reader to notice.

Contracting the tensor result with $\eta_{\mu\nu}$ agrees with
\[
  \int\frac{\dd^{d_{\mathrm{reg}}}p}{(2\pi)^{d_{\mathrm{reg}}}}
  \frac{p^2}{(p^2+U-\ii0)^n}
  =I_{n-1}(U)-U I_n(U)
\]
after using $\Gamma(z+1)=z\Gamma(z)$.  The contraction and the dimensions
check the factor $1/2$, the power of $U$, and the explicit
$\eta^{\mu\nu}$.

\BanksImplicitSolution{I-CH09-022}
Near a pole,
\[
  \mathcal D\left(\frac{Z_i}{p^2+M_i^2-\ii0}\right)
  =\frac{\mathcal D Z_i}{p^2+M_i^2-\ii0}
  +\frac{Z_i\mathcal D M_i^2}
  {(p^2+M_i^2-\ii0)^2}.
\]
The RG equation has only the original simple pole on its right-hand side.
The double-pole coefficient must therefore vanish:
\[
  \mathcal D M_i^2=0.
\]
Equating the simple-pole residues gives
\[
  \mathcal D Z_i=-2\gamma_iZ_i.
\]
The separation works because $\mathcal D$ is first order.  A second-order
operator would also differentiate products of parameter-dependent factors and
would spoil this simple residue argument.

For an amplitude with one LSZ factor per external leg, write
\[
  S=\lim_{\mathrm{on\ shell}}
  \left[\prod_i f_i(p_i)(p_i^2+M_i^2)Z_i^{-1/2}\right]W_N.
\]
The normalized wave functions depend on RG parameters through physical masses,
so $\mathcal Df_i=0$.  We also have
\[
  \mathcal D(p_i^2+M_i^2)=0,\qquad
  \mathcal D Z_i^{-1/2}=+\gamma_iZ_i^{-1/2},
\]
while
\[
  \mathcal DW_N=-\left(\sum_i\gamma_i\right)W_N.
\]
The first-order Leibniz rule now gives
\begin{align*}
  \mathcal DS
  &=\left(\sum_i\gamma_i\right)S
  -\left(\sum_i\gamma_i\right)S=0.
\end{align*}
The external-leg anomalous dimensions cancel the Green-function anomalous
dimensions one by one.

This derivation assumes isolated stable one-particle poles, a mass-independent
RG operation on the external momentum labels, and infrared control sufficient
to interchange $\mathcal D$ with the LSZ limit.  Charged states in QED require
the dressed or inclusive refinement discussed by Banks.  In a free theory,
$\beta=\gamma_i=0$, $Z_i=1$, and the equation reduces immediately to
$\mu\partial_\mu S=0$.

\BanksImplicitSolution{I-CH09-023}
For a graph with $I$ internal NGB lines, $V$ vertices, and $L$ loops, loop
measures supply $p^{4L}$, propagators supply $p^{-2I}$, and a vertex $v$
supplies $p^{d_v}$.  Its momentum degree is
\[
  \omega=4L-2I+\sum_vd_v.
\]
For a connected graph, $L=I-V+1$.  Substitution gives
\[
  \omega=2+2L+\sum_v(d_v-2).
\]
Every loop or replacement of a leading two-derivative vertex by a
higher-derivative vertex raises the low-momentum order.

Expanding $U=\exp(\ii\pi^aT^a/f)$ shows that a leading $n$-field vertex scales
as $p^2/f^{n-2}$.  Combining all vertices and using the line-counting identity
$\sum_vn_v=2I+E$ leaves the overall factor $f^{2-E}$.  Each loop integration
also carries $(4\pi)^{-2}$ in four dimensions.  With natural dimensionless
low-energy constants, the result can be organized as
\[
  \mathcal A_E\sim\frac{p^2}{f^{E-2}}
  \left(\frac{p}{4\pi f}\right)^{2L}
  \prod_v\left(\frac{p}{f}\right)^{d_v-2}.
\]
Group tensors and angular functions multiply this estimate without changing
its momentum degree.

Thus a tree graph from the two-derivative action starts at
$p^2/f^{E-2}$.  A four-derivative insertion is down by $(p/f)^2$, while one
loop is down by $(p/4\pi f)^2$.  More derivatives and more loops generate a
graded expansion.  Naive dimensional analysis places its cutoff near
$4\pi f$; the conservative regime $p\ll f$ used in Banks's sentence makes
every displayed ratio small.

Massless loops generate nonanalytic terms such as
$p^4\ln(-p^2)/[(4\pi)^2f^4]$.  Their analytic polynomial pieces have the same
form as local higher-derivative operators and are absorbed into the
renormalized low-energy constants.  The nonlinear-symmetry Ward identity gives
the Adler zero when one external NGB momentum is taken soft, after all diagrams
at a fixed order are summed.  Derivative counting fixes the overall power of
momenta; it cannot prove that zero diagram by diagram because an expanded
vertex may contain undifferentiated pion fields.  The cancellations among the
tree graphs of the leading nonlinear sigma model provide the first check.
Sending $p/f\to0$ removes every higher-order correction and leaves the leading
action.

\end{bankssolutions}
